\documentclass{article}
\usepackage[margin=2cm]{geometry}
\usepackage{amsmath}
\usepackage{graphicx}
\usepackage{booktabs}
\usepackage[utf8]{inputenc}
\usepackage{ragged2e}
\setlength{\emergencystretch}{3em}
\usepackage{paracol}
\begin{document}
\sloppy

\noindent\hfill\includegraphics[height=1.8em]{assets/figure_header_logo.png}\par\noindent\hrule\vspace{4pt}

\begin{paracol}{2}
\RaggedRight

with implied star product and symmetric C components. Master field components are now separated in
physical scalar field C and higher ones, related to it onshell by derivatives. The expansion of C is given by


\begin{equation}
C = \sum _ { s = 1 } ^ { \infin } \sum _ { | m | < s } C _ { m } ^ { s } V _ { m } ^ { s } \qquad ( 1 )
\end{equation}


g1 versus j2 are related with 2m = N - N2 and Cm are
functions of spacetime coordinates. Auxiliary tensors are
absorbed within a field. The fields A and A are expanded
analogously


\begin{equation}
A = \sum _ { s = 2 } ^ { \infin } \sum _ { | m | < s } A _ { m } ^ { s } V _ { m } ^ { s } ~ ~ ~ ~ ~ ~ \overline { A } = \sum _ { s = 2 } ^ { \infin } \sum _ { | m | < s } \overline { A } _ { m } ^ { s } V _ { m } ^ { s }. ~ ~ ~ ~ ~ ( 2 )
\end{equation}


The standard procedure of finding the generalised KG
equation consists of inserting the expressions for A, A
and C in O and determining the smallest possible set
of equations needed to find the scalar equation in arbitrary background. Standard procedure can be described
considering equation O in AdS background since it is a
foundation for the following computations. The vacuum
C equation without AdS fields is ordinary KG equation
while one can determine the higher components in the
terms of Co. The AdS connection consists of the spin-2
generators that form SL(2), subalgebra of hs[X]


\begin{equation}
\begin{array} { l } { A = e ^ { \rho } V _ { 1 } ^ { 2 } d z + V _ { 0 } ^ { 2 } d \rho \qquad ( 3 ) } \\ { \bar { A } = e ^ { \rho } V _ { - 1 } ^ { 2 } d \bar { z } - V _ { 0 } ^ { 2 } d \rho \qquad \qquad ( 4 ) } \end{array}
\end{equation}


with AdS metric


The higher spins felds vanish, and we are working in Euclidean metric and Fefferman-Graham gauge. The general form of the C equation O in the AdS background


\begin{equation}
\partial _ { \rho } C _ { m } ^ { s } + 2 C _ { m } ^ { s + 1 } + C _ { m } ^ { s + 1 } g _ { 3 } ^ { ( s + 1 ) 2 } ( m, 0 ) = 0 \qquad ( 6 )
\end{equation}


\begin{equation}
\begin{array} { r } { \partial C _ { m } ^ { s } + e ^ { \rho } ( C _ { m - 1 } ^ { s - 1 } + \frac { 1 } { 2 } g _ { 2 } ^ { 2 s } ( 1, m - 1 ) C _ { m } ^ { s } } \\ { + \frac { 1 } { 2 } g _ { 3 } ^ { 2 ( s + 1 ) } ( 1, m - 1 ) C _ { m - 1 } ^ { s + 1 } ) = 0 } \end{array}
\end{equation}


\begin{equation}
\begin{array} { c } { \overline { \partial } C _ { m } ^ { s } - e ^ { \rho } ( C _ { m + 1 } ^ { s - 1 } - \frac { 1 } { 2 } g _ { 2 } ^ { 2 s } ( - 1, m + 1 ) C _ { m + 1 } ^ { s } } \\ { + \frac { 1 } { 2 } g _ { 3 } ^ { 2 ( s + 1 ) } ( - 1, m + 1 ) C _ { m + 1 } ^ { s + 1 } ) = 0 } \end{array}
\end{equation}


for m < s, 0 = z,  = = and the X-dependence in
the structure constants suppressed. In the simplest case
choosing s = 1, s = 2 one can solve for the higher components in C and obtain the Klein-Gordon KG equation


\begin{equation}
\left [ \partial _ { \rho } ^ { 2 } + 2 \partial _ { \rho } + 4 e ^ { - 2 \rho } \partial \bar { \partial } - ( \lambda ^ { 2 } - 1 ) \right ] C _ { 0 } ^ { 1 } = 0. \qquad ( 9 )
\end{equation}

\switchcolumn
\RaggedRight

Consistency condition on equations is that all the components of C have smooth solution when expressed using
Cl. The strategy for determining the minimal set of
equations for C is to select components of C that are of
 m and therefore the smallest spin for fixed
 equations for fixed m, solve for non-minimal
components in terms of minimal ones and p derivatives,
for Ap = -A, = V2. After solving for minimal ones,
Cl and its derivatives. Once that we have expressed the
higher components of C in terms of C we can determine
the part that defines the KG equation and the generalised
part that appears due to the HS background. To obtain
the equation of motion for the scalar field up to linear
order we consider the variation of the gauge field and
apply the KG equation on it. This and the standard procedure for obtaining the linearised equation of motion for
the scalar field described above should be equal once the
gauge parameter is chosen conveniently. That approach
can be written in the following way. First we express the
higher components of the C field in terms of the combination of the derivatives on C in the background AdS.
Focusing on the master field C, the equation O is invariant under the hs[X]   hs[X] gauge invariance when


\begin{equation}
C \rightarrow C + C \star \bar { \Lambda } - \Lambda \star C \qquad \qquad ( 1 0 )
\end{equation}


\begin{equation}
\Lambda ( \rho, z, \bar { z } ) = \sum _ { n = 1 } ^ { 2 s - 1 } \frac { 1 } { ( n - 1 )! } ( - \partial ) ^ { n - 1 } \lambda ^ { ( s ) } ( z, \bar { z } ) e ^ { ( s - n ) \rho } V _ { s - n } ^ { s }.
\end{equation}


Where we take A to be chiral, so A = O. The feld in the
higher spin background is obtained by transformation


\begin{equation}
\tilde { C } _ { m } ^ { s } = C _ { m } ^ { s } - ( \Lambda \star C ) _ { m } ^ { s }. \qquad \qquad ( 1 2 )
\end{equation}


The field Cm we express in terms of the Cl. To do that
we focus on the set of equations O,(20),O). The product
of the C field with A gives combination of higher components of C in AdS background which can, as we will
show, be expressed in terms of C. On the field Cl we
can use the transformation O and obtain


\begin{equation}
\tilde { C } _ { 0 } ^ { 1 } = C _ { 0 } ^ { 1 } - ( \Lambda \star C ) _ { 0 } ^ { 1 }. \qquad \qquad ( 1 3 )
\end{equation}


Since we are at the linear order, once we have C we
can rewrite it as Cl which is defined on the higher spin
background. From the expression for the gauge field A O
and the relation for the star product O we can determine
the variation of the scalar field C


\begin{equation}
( \delta C ) _ { 0 } ^ { 1 } = - \sum _ { n = 1 } ^ { 2 s - 1 } \frac { 1 } { ( n - 1 )! } ( - \partial ) ^ { n - 1 } \Lambda ^ { ( s ) } \qquad \qquad ( 1 4 )
\end{equation}


\begin{equation}
\times ~ \frac { 1 } { 2 } g _ { 2 s - 1 } ^ { s s } ( s - n, n - s ) C _ { - ( s - n ) } ^ { s } e ^ { ( s - n ) \rho }, ~ ~ ~ ~ ( 1 5 )
\end{equation}

\end{paracol}

\end{document}
